\documentclass[12pt]{article}
\usepackage[utf8]{inputenc}
\usepackage[T1]{fontenc}
\usepackage{hyperref}
\title{Sudovian Language Resources for K-12 Curriculum Development}
\author{Compiled by Claude Sonnet 5.5 (claude-sonnet-5-5) for Andr\'as Kornai}
\date{}
\begin{document}
\maketitle

\section{Introduction}
Sudovian (Yotvingian, Jatvingian; ISO 639-3: xsv) is an extinct West Baltic language. It is attested only in very small written remains, and scholars dispute whether it was separate from Old Prussian. This report rests on one web search. It found background and no teaching materials.

\section{Primary Educational Resources}
No primers, textbooks or school programmes were found. Any heritage curriculum would be about the people, history and attested words, not a living language.

\section{Linguistic and Reference Materials}
\begin{itemize}
\item Wikipedia, Sudovian language: overview and the debate on its classification. \url{https://en.wikipedia.org/wiki/Sudovian_language}
\item Wikipedia, Yotvingians: the people, and modern revival movements. \url{https://en.wikipedia.org/wiki/Yotvingians}
\item ``Jotvingi\c{u} kalbos aspektai'', Lituanistika (Lithuanian-language scholarship). Not examined in detail. \url{https://www.lituanistika.lt/content/59440}
\item ``Seniausios \v{z}inios apie s\={u}duvius-jotvingius ir j\c{u} \.{i}vardijimo problema'', Lituanistika. \url{https://lituanistika.lt/content/104730}
\end{itemize}

\section{Religious and Community Resources}
The Ukrainian Week article ``Terra Sudorum'' is a popular history piece and mentions revival attempts and the fragments from which recreation is attempted. It needs verification against scholarly sources. \url{https://ukrainianweek.com/terra-sudorum/}

\section{Digital and Audio Resources}
The search found none. The Lexvo entry for xsv is only a language-identification record. \url{https://finto.fi/lexvo/en/page/xsv}

\section{Curriculum Development Context}
Little exists. The search found no active revival programme and no Polish, Belarusian or Lithuanian school curriculum. A heritage unit could draw on the scholarly articles and the Wikipedia overviews, treating any reconstruction as hypothetical.
\end{document}